\begin{resumo}
 A qualidade de produtos de software é essencial para garantir a 
 manutenibilidade e reduzir custos ao longo do ciclo de vida. Modelos de 
 qualidade como ISO/IEC 25010 destacam a manutenibilidade como característica 
 fundamental, composta por atributos como analisabilidade, modificabilidade e 
 testabilidade. Os code smells, indícios de problemas de projeto e 
 implementação, estão associados a maior esforço de manutenção e degradação 
 da qualidade interna. Ferramentas tradicionais de análise estática, como 
 SonarQube, são amplamente utilizadas na detecção desses problemas, porém 
 apresentam limitações relacionadas a falsos positivos e dependência de 
 limiares fixos de métricas. Recentemente, técnicas de Inteligência Artificial 
 (IA), especialmente Modelos de Linguagem de Grande Escala (LLMs), 
 têm demonstrado potencial para capturar padrões estruturais e semânticos 
 mais complexos. Este trabalho investiga empiricamente como sistemas baseados 
 em IA auxiliam a garantia de qualidade de software na detecção de anomalias 
 em código-fonte. Para isso, foi realizada uma revisão estruturada da literatura, 
 adaptada do protocolo de revisão  sistemática, que identificou 42 estudos relevantes 
 publicados entre 2022 e 2025. A análise evidenciou que, embora ferramentas de 
 análise estática continuem predominantes na indústria, abordagens baseadas em IA 
 apresentam capacidade de identificação de padrões que extrapolam 
 regras sintáticas fixas. Foi então planejado um experimento controlado que compara as 
 LLMs com o SonarQube na detecção de code smells, utilizando datasets anotados
 manualmente como referência. Os resultados esperados visam produzir evidências empíricas 
 sobre a contribuição relativa de cada abordagem, identificando limitações e oportunidades 
 de integração entre métodos tradicionais e técnicas emergentes de IA na garantia de 
 qualidade de código.

 \vspace{\onelineskip}
    
 \noindent
 \textbf{Palavras-chave}: engenharia de software. inteligência artificial. code smells. qualidade de software. modelos de linguagem. análise estática.
\end{resumo}
